% answers_short.tex —— 1985 年·广东卷·理工农医类 · 简短答案速查
% 由 tools/make_short_answers.py 自 answers.tex 的 \gksolhead/\gkanswer 提取，
% 并与 1985/1985_数学真题_广东理工农医类_答案要点.md 逐题比对后生成，勿手改（改源 MD 后重跑）。

\begin{gknotice}
\textbf{使用说明}：本节只给最终答案便于快速核对，完整推导见解析版。答案为编写者独立推导，\textbf{非官方标准答案}。
\end{gknotice}

\gkpart{第一卷填空与画图}
\begin{itemize}[leftmargin=2.6em]
  \item[\textbf{1.}] $x\leq1$。
  \item[\textbf{2.}] $\dfrac{x}{x^2-1}$（题设只确定 $x\ne0,\pm1$ 时的函数值；$f(0)$ 未定）。
  \item[\textbf{3.}] $\dfrac19$。
  \item[\textbf{4.}] $x^2+y^2-x-2y=0$。
  \item[\textbf{5.}] $0$。
  \item[\textbf{6.}] $\dfrac32$。
  \item[\textbf{7.}] $\dfrac\pi2$。
  \item[\textbf{8.}] $-10$。
  \item[\textbf{9.}] $y=2+\sqrt{9-x^2}$，$-3\leq x\leq3$。
  \item[\textbf{10.}] 曲线为 $x=1$ 上的整条直线与 $y=2\sqrt x\ (x\geq0)$ 的并集。
\end{itemize}

\gkpart{第一卷选择}
\begin{itemize}[leftmargin=2.6em]
  \item[\textbf{1.}] E。
  \item[\textbf{2.}] A。
  \item[\textbf{3.}] B。
  \item[\textbf{4.}] A。
  \item[\textbf{5.}] C。
  \item[\textbf{6.}] C。
  \item[\textbf{7.}] E。
  \item[\textbf{8.}] B。
  \item[\textbf{9.}] D。
  \item[\textbf{10.}] C。
  \item[\textbf{11.}] C。
  \item[\textbf{12.}] B。
  \item[\textbf{13.}] D。
  \item[\textbf{14.}] C。
  \item[\textbf{15.}] D。
  \item[\textbf{16.}] D。
  \item[\textbf{17.}] C。
  \item[\textbf{18.}] A。
  \item[\textbf{19.}] E。
  \item[\textbf{20.}] B。
\end{itemize}

\gkpart{第二卷大题}
\begin{itemize}[leftmargin=2.6em]
  \item[\textbf{一.}] 见解析版
  \item[\textbf{二.}] (1) $CE\perp AB_1$；(2) $\dfrac{2\sqrt3}{3}\,\mathrm{cm}$；(3) $\cos\alpha=\dfrac{\sqrt{10}}5$
  \item[\textbf{三.}] (1) $m>1$；(2) $\displaystyle\min_x f(x)=\log_3\!\left(m+\frac1{m-1}\right)$，在 $x=2m$ 取到；(3) $\min_x f(x)\geq1$
  \item[\textbf{四.}] (1) 当 $m=4$ 时，共同焦点为 $(3,0)$、共同准线为 $x=0$；(2) 两曲线有交点当且仅当 $m\in[-\tfrac12,\tfrac72]$；(3) $m=0$ 或 $m=8(\sqrt2-1)$
  \item[\textbf{五.}] (1) $x=-1$ 时极大值为 $c+4$；$x=1$ 时极小值为 $c-4$；(2) 见解析版
\end{itemize}
